\par\noindent\textbf{(1)} 物理场绕轴旋转φ→φ+2π后必须单值，要求exp(i2πℓ)=1，因此ℓ∈整数。相位梯度∇\_⊥(ℓφ)=ℓ eφ/r，其绕任何包围光轴的闭环积分为∮∇Φ·dl=2πℓ，称相位拓扑荷。
\par\noindent\textbf{(2)} 局部横向波矢的方位分量k\_φ=ℓ/r，因而傍轴能流可有绕轴的角向分量。ℓ>0与ℓ<0的涡旋流环绕方向相反；r→0处k\_φ形式上发散，要求场幅F趋零来使光学场和能流保持正则。
\par\noindent\textbf{(3)} 关于z轴的轨道角动量算符L\_z=−iℏ∂\_φ。作用于U=Fexp(iℓφ)得L\_z U=ℓℏ U，因此理想纯ℓ模式每个光子的z向轨道角动量为ℓℏ，若混合多模式则平均值是模式权重加权而非单个整数。
\par\noindent\textbf{(4)} 正则ℓ≠0涡旋的近轴振幅通常F(r)∝r\textasciicircum{}\{|ℓ|\}，强度I∝r\textasciicircum{}\{2|ℓ|\}，中心为暗核；拓扑荷反号改变绕轴相位方向而不一定改变径向强度。强非傍轴场会有自旋—轨道混合，不能把所有角动量无条件归为标量ℓℏ。
